\documentclass[9pt,a4paper,twocolumn,twoside]{rho-class/rho}

\usepackage[english]{babel}
\usepackage[utf8]{inputenc}
\usepackage[T1]{fontenc}
\usepackage{booktabs}
\usepackage{algorithm}
\usepackage{algpseudocode}
\usepackage{pgfplots}
\usepackage{tikz}
\usepackage{float}
\usepackage{comment}
\usepackage{hyperref}
\usepackage{array}

\usetikzlibrary{shapes.geometric, arrows.meta, positioning}
\tikzset{
    flow/.style={node distance=0.75cm, every node/.style={font=\footnotesize}},
    startstop/.style={rectangle, rounded corners, minimum width=1.8cm, minimum height=0.5cm, draw, align=center},
    process/.style={rectangle, minimum width=2.2cm, minimum height=0.5cm, draw, align=center},
    decision/.style={diamond, aspect=2, draw, align=center, inner sep=1pt, text width=2.0cm},
    arrow/.style={thick, -{Latex[length=2mm]}}
}

\doctype{Reporte de Diseño}
\title{Documentación de decisiones de diseño para asignación de pines y control de alimentación en "Placa Madre Mini"}

\author[{\textasteriskcentered},a]{Adonay, Sebastián}
\affil[a]{Centro de Investigación en Tecnologías para la Sociedad C+, Universidad del Desarrollo}
\dates{Este manuscrito ha sido compilado el \today}

\corres{\textsuperscript{\textasteriskcentered}Autor correspondiente: \href{mailto:s.adonay@udd.cl}{s.adonay@udd.cl} (S. Adonay).}

\docinfo{Documento interno de apoyo al diseño electrónico y firmware.}

\journalname{Documentación Interna}
\journal{C+ UDD 2026}
\theday{\today}

\begin{abstract}
Se documentan las principales decisiones de diseño adoptadas para la asignación de pines de una placa basada en ESP32-S3-WROOM-1, incluyendo buses de comunicación, expansión mediante conector IDC, medición de batería, control de NeoPixel y manejo de rieles de alimentación conmutables para periféricos y módem. Además, se explicita la justificación para reservar ciertos GPIO y se resume el pinout global e interfaz de expansión.
\end{abstract}

\keywords{ESP32-S3, Pinout, GPIO, I2C, SPI, UART, NeoPixel, Power Gating, SIM7600, IDC}

\begin{document}

\maketitle
\thispagestyle{firststyle}

\section{Introducción}

\rhostart{E}l presente documento resume las decisiones de diseño adoptadas para una placa basada en \textbf{ESP32-S3-WROOM-1}, con énfasis en la asignación funcional de pines, la segregación de rieles de alimentación y la exposición de señales hacia un conector de expansión IDC.

El objetivo principal fue construir una arquitectura ordenada y mantenible, adecuada tanto para desarrollo de firmware como para revisión de hardware, evitando conflictos con pines sensibles de arranque, USB nativo y señales con restricciones funcionales.

\section{Descripción general del sistema}

La arquitectura considera tres rieles principales de interés:

\begin{itemize}
    \item \texttt{3V3}: riel principal siempre activo.
    \item \texttt{3V3\_INT}: riel 3.3V conmutable para periféricos generales y externos.
    \item \texttt{3V3\_MODEM}: riel 3.3V conmutable dedicado al módem.
    \item \texttt{5V\_INT}: riel 5.0V conmutable para sensor "Plantower" y periféricos externos.
\end{itemize}

El microcontrolador mantiene alimentación permanente, mientras que periféricos de mayor consumo o de operación intermitente se energizan selectivamente mediante MOSFET P-channel en configuración \textit{high-side}.

\section{Decisiones de diseño}

\subsection{Buses por defecto}

Se decidió mantener los buses principales usando la asignación estándar del entorno Arduino IDE para la placa \texttt{ESP32S3 Dev Module}, con el fin de facilitar compatibilidad, depuración y legibilidad del firmware:

\begin{itemize}
    \item \textbf{I2C}: GPIO8 (SDA) y GPIO9 (SCL).
    \item \textbf{SPI}: GPIO10 (SS), GPIO11 (MOSI), GPIO12 (SCK) y GPIO13 (MISO).
\end{itemize}

En el caso particular del conector IDC, la señal \texttt{SS} no fue exportada, debido a que la tarjeta SD reside en la placa principal y utiliza localmente dicha línea.

\newpage

\subsection{UARTs}

Se reservaron UARTs diferenciadas para propósitos distintos:

\begin{itemize}
    \item UART auxiliar en GPIO4 y GPIO5, conectado a conector "Plantower"
    \item UART dedicada al módem en GPIO15 y GPIO16.
    \item UART0 de programación y depuración en GPIO43 y GPIO44.
\end{itemize}

\subsection{Medición de batería}

La medición analógica del voltaje de batería se fijó en \textbf{GPIO1}, correspondiente a un canal ADC válido del ESP32-S3.

\subsection{NeoPixel}

La señal de control del NeoPixel se asignó a \textbf{GPIO14}. Este pin también se mantuvo disponible en el conector IDC, permitiendo flexibilidad de uso futuro.

\subsection{Power gating}

Se adoptó control de rieles mediante MOSFET P-channel:\\
\textbf{Nota: los controles de gate 3v3 y modem se activan en LOW y el de 5v en HIGH}
\begin{itemize}
    \item \textbf{GPIO41}: control de \texttt{Gate\_3V3INT}.
    \item \textbf{GPIO42}: control de \texttt{Gate\_MODEM}.
\end{itemize}

La conmutación se realiza por el lado alto (\textit{high-side}), evitando corte de GND y reduciendo riesgos de retroalimentación a través de interfaces lógicas.

A su vez, el riel de 5V se activa al entregar voltaje al pin ENABLE del regulador AP63357 de 5V.
\begin{itemize}
    \item \textbf{GPIO2}: control de \texttt{EN\_5V\_INT}.
\end{itemize}


\subsection{Pines reservados o evitados}

Se decidió evitar o reservar los siguientes GPIO por su naturaleza funcional:

\begin{itemize}
    \item GPIO0 y GPIO3: pines de \textit{strapping}.
    \item GPIO19 y GPIO20: USB nativo.
    \item GPIO45: pin de \textit{strapping}.
    \item GPIO46: pin de entrada solamente, además de \textit{strapping}.
\end{itemize}




\newpage

\section{Asignación por categoría}

\begin{table}[H]
\centering
\resizebox{0.55\columnwidth}{!}{
\begin{tabular}{ccl}
\toprule
\textbf{GPIO} & \textbf{Categoría} & \textbf{Función} \\
\midrule
\multicolumn{3}{l}{\textbf{I2C}} \\
\midrule
8  & I2C   & SDA \\
9  & I2C   & SCL \\

\midrule
\multicolumn{3}{l}{\textbf{SPI / SD}} \\
\midrule
10 & SPI   & SS (SD en placa principal) \\
11 & SPI   & MOSI \\
12 & SPI   & SCK \\
13 & SPI   & MISO \\

\midrule
\multicolumn{3}{l}{\textbf{UART}} \\
\midrule
4  & UART  & RX auxiliar \\
5  & UART  & TX auxiliar \\
15 & UART  & RX módem \\
16 & UART  & TX módem \\
43 & UART  & TX0 (debug) \\
44 & UART  & RX0 (debug) \\

\midrule
\multicolumn{3}{l}{\textbf{ADC}} \\
\midrule
1  & ADC   & Medición de batería \\

\midrule
\multicolumn{3}{l}{\textbf{Power Control}} \\
\midrule
2 & Power & \texttt{EN\_5V\_INT} \\
41 & Power & \texttt{Gate\_3V3\_INT} \\
42 & Power & \texttt{Gate\_MODEM} \\

\midrule
\multicolumn{3}{l}{\textbf{NeoPixel}} \\
\midrule
14 & LED   & Data NeoPixel \\

\midrule
\multicolumn{3}{l}{\textbf{Expansión (IDC)}} \\
\midrule
14 & IDC   & GPIO expansión \\
21 & IDC   & GPIO expansión \\
35 & IDC   & GPIO expansión \\
36 & IDC   & GPIO expansión \\
37 & IDC   & GPIO expansión \\
38 & IDC   & GPIO expansión \\
39 & IDC   & GPIO expansión \\
40 & IDC   & GPIO expansión \\
47 & IDC   & GPIO expansión \\
48 & IDC   & GPIO expansión \\

\midrule
\multicolumn{3}{l}{\textbf{Reservados / Evitados}} \\
\midrule
0  & Reserved & Strapping (boot) \\
3  & Reserved & Strapping (boot) \\
19 & Reserved & USB D- \\
20 & Reserved & USB D+ \\
45 & Reserved & Strapping \\
46 & Reserved & Input-only + strapping \\
\bottomrule
\end{tabular}
}
\caption{Asignación de pines agrupada por categoría}
\end{table}

\section{Asignación por Pin GPIO ESP32}

\begin{table}[H]
\centering
\resizebox{0.7\columnwidth}{!}{
\begin{tabular}{cll}
\toprule
\textbf{GPIO} & \textbf{Categoría} & \textbf{Función} \\
\midrule
0   & Strapping & No usar (boot)\\
1   & ADC & Lectura batería (ADC1\_CH0) \\
2   & Power Control & Pin Enable de StepDown $5V_{INT}$\\
3   & Strapping & No usar (boot)\\
4   & UART & RX auxiliar \\
5   & UART & TX auxiliar \\
6   & GPIO & Botón 1 \\
7   & GPIO & Botón 2 \\
8   & I2C & SDA \\
9   & I2C & SCL \\
10  & SPI & SS (SD en placa principal) \\
11  & SPI & MOSI \\
12  & SPI & SCK \\
13  & SPI & MISO \\
14  & GPIO / LED & NeoPixel (también exportado a IDC) \\
15  & UART & RX módem \\
16  & UART & TX módem \\
17  & GPIO (JST) & Expansión \\
18  & GPIO (JST) & Expansión \\
19  & USB D- & Pin D- de comunicación USB \\
20  & USB D+ & Pin D+ de comunicación USB \\
21  & GPIO (IDC) & Expansión \\
35  & GPIO (IDC) & Expansión \\
36  & GPIO (IDC) & Expansión \\
37  & GPIO (IDC) & Expansión \\
38  & GPIO (IDC) & Expansión \\
39  & GPIO (IDC) & Expansión \\
40  & GPIO (IDC) & Expansión \\
41  & Power Control & \texttt{Gate\_3V3INT} \\
42  & Power Control & \texttt{Gate\_MODEM} \\
43  & UART & TX0 (programación / debug) \\
44  & UART & RX0 (programación / debug) \\
45  & Strapping & No usar \\
46  & Strapping & No usar \\
47  & GPIO (IDC) & Expansión \\
48  & GPIO (IDC) & Expansión \\
\bottomrule
\end{tabular}
}
\caption{Asignación global de pines del sistema}
\label{tab:pinout_global}
\end{table}



\newpage%-------------------------------------------------\section{Asignación conectores}

\begin{table}[H]
\centering
\setlength{\tabcolsep}{4pt}
\begin{tabular}{cll}
\toprule
\textbf{Pin} & \textbf{Señal} & \textbf{Descripción} \\
\midrule

\multicolumn{3}{l}{\textbf{IDC-26 (J8)}} \\
\midrule
1  & \texttt{12V}     & Alimentación 12V \\
2  & \texttt{12V}     & Alimentación 12V \\
3  & \texttt{5V\_INT} & Alimentación 5V \\
4  & \texttt{5V\_INT} & Alimentación 5V \\
5  & \texttt{3V3\_INT} & Rail conmutable \\
6  & \texttt{3V3\_INT} & Rail conmutable \\
7  & \texttt{3V3}      & Rail principal \\
8  & \texttt{IO47}     & GPIO \\
9  & \texttt{IO21}     & GPIO \\
10 & \texttt{IO48}     & GPIO \\
11 & \texttt{IO14}     & GPIO / NeoPixel \\
12 & \texttt{IO35}     & GPIO \\
13 & \texttt{IO13}     & SPI MISO \\
14 & \texttt{IO36}     & GPIO \\
15 & \texttt{IO12}     & SPI SCK \\
16 & \texttt{IO37}     & GPIO \\
17 & \texttt{IO11}     & SPI MOSI \\
18 & \texttt{IO38}     & GPIO \\
19 & \texttt{IO9}      & I2C SCL \\
20 & \texttt{IO39}     & GPIO \\
21 & \texttt{IO8}      & I2C SDA \\
22 & \texttt{IO40}     & GPIO \\
23 & \texttt{GND}      & Tierra \\
24 & \texttt{GND}      & Tierra \\
25 & \texttt{GND}      & Tierra \\
26 & \texttt{GND}      & Tierra \\

\midrule
\multicolumn{3}{l}{\textbf{Plantower (J1)}} \\
\midrule
1 & \texttt{N/C} & --- \\
2 & \texttt{N/C} & --- \\
3 & \texttt{N/C} & --- \\
4 & \texttt{IO4} & UART RX \\
5 & \texttt{IO5} & UART TX \\
6 & \texttt{N/C} & --- \\
7 & \texttt{GND} & Tierra \\
8 & \texttt{5V\_INT} & Alimentación \\

\midrule
\multicolumn{3}{l}{\textbf{I2C (J2)}} \\
\midrule
1 & \texttt{GND} & Tierra \\
2 & \texttt{3V3\_INT} & Alimentación \\
3 & \texttt{IO9} & SCL \\
4 & \texttt{IO8} & SDA \\

\midrule
\multicolumn{3}{l}{\textbf{I2C (J3)}} \\
\midrule
1 & \texttt{GND} & Tierra \\
2 & \texttt{3V3\_INT} & Alimentación \\
3 & \texttt{IO9} & SCL \\
4 & \texttt{IO8} & SDA \\

\midrule
\multicolumn{3}{l}{\textbf{Botón (J4)}} \\
\midrule
1 & \texttt{GND} & Tierra \\
2 & \texttt{IO6} & Entrada botón \\

\midrule
\multicolumn{3}{l}{\textbf{Botón (J5)}} \\
\midrule
1 & \texttt{GND} & Tierra \\
2 & \texttt{IO7} & Entrada botón \\

\midrule
\multicolumn{3}{l}{\textbf{GPIO (J6)}} \\
\midrule
1 & \texttt{GND} & Tierra \\
2 & \texttt{3V3\_INT} & Alimentación \\
3 & \texttt{IO18} & GPIO \\
4 & \texttt{IO17} & GPIO \\

\midrule
\multicolumn{3}{l}{\textbf{NeoPixel (J7)}} \\
\midrule
1 & \texttt{GND} & Tierra \\
2 & \texttt{3V3} & Alimentación \\
3 & \texttt{IO14} & Data \\

\bottomrule
\end{tabular}
\caption{Resumen completo de conectores y pinout}
\end{table}

\printbibliography

\end{document}